% Options for packages loaded elsewhere
\PassOptionsToPackage{unicode}{hyperref}
\PassOptionsToPackage{hyphens}{url}
%
\documentclass[
  12pt,
  a4paper,
  DIV=11,
  numbers=noendperiod]{scrreprt}

\usepackage{amsmath,amssymb}
\usepackage{setspace}
\usepackage{iftex}
\ifPDFTeX
  \usepackage[T1]{fontenc}
  \usepackage[utf8]{inputenc}
  \usepackage{textcomp} % provide euro and other symbols
\else % if luatex or xetex
  \usepackage{unicode-math}
  \defaultfontfeatures{Scale=MatchLowercase}
  \defaultfontfeatures[\rmfamily]{Ligatures=TeX,Scale=1}
\fi
\usepackage{lmodern}
\ifPDFTeX\else  
    % xetex/luatex font selection
\fi
% Use upquote if available, for straight quotes in verbatim environments
\IfFileExists{upquote.sty}{\usepackage{upquote}}{}
\IfFileExists{microtype.sty}{% use microtype if available
  \usepackage[]{microtype}
  \UseMicrotypeSet[protrusion]{basicmath} % disable protrusion for tt fonts
}{}
\usepackage{xcolor}
\setlength{\emergencystretch}{3em} % prevent overfull lines
\setcounter{secnumdepth}{5}
% Make \paragraph and \subparagraph free-standing
\ifx\paragraph\undefined\else
  \let\oldparagraph\paragraph
  \renewcommand{\paragraph}[1]{\oldparagraph{#1}\mbox{}}
\fi
\ifx\subparagraph\undefined\else
  \let\oldsubparagraph\subparagraph
  \renewcommand{\subparagraph}[1]{\oldsubparagraph{#1}\mbox{}}
\fi


\providecommand{\tightlist}{%
  \setlength{\itemsep}{0pt}\setlength{\parskip}{0pt}}\usepackage{longtable,booktabs,array}
\usepackage{calc} % for calculating minipage widths
% Correct order of tables after \paragraph or \subparagraph
\usepackage{etoolbox}
\makeatletter
\patchcmd\longtable{\par}{\if@noskipsec\mbox{}\fi\par}{}{}
\makeatother
% Allow footnotes in longtable head/foot
\IfFileExists{footnotehyper.sty}{\usepackage{footnotehyper}}{\usepackage{footnote}}
\makesavenoteenv{longtable}
\usepackage{graphicx}
\makeatletter
\def\maxwidth{\ifdim\Gin@nat@width>\linewidth\linewidth\else\Gin@nat@width\fi}
\def\maxheight{\ifdim\Gin@nat@height>\textheight\textheight\else\Gin@nat@height\fi}
\makeatother
% Scale images if necessary, so that they will not overflow the page
% margins by default, and it is still possible to overwrite the defaults
% using explicit options in \includegraphics[width, height, ...]{}
\setkeys{Gin}{width=\maxwidth,height=\maxheight,keepaspectratio}
% Set default figure placement to htbp
\makeatletter
\def\fps@figure{htbp}
\makeatother

\usepackage{indentfirst}
\usepackage{float}
\floatplacement{figure}{H}
\IfFileExists{plex-otf.sty}{
  %% Full TeXlive
  \usepackage[
    % math,
    RM={Scale=0.94},SS={Scale=0.94},SScon={Scale=0.94},TT={Scale=MatchLowercase,FakeStretch=0.9},
    DefaultFeatures={Ligatures=Common}
  ]{plex-otf}
  % \usepackage[Scale=MatchUppercase]{juliamono}
}{
  %% TinyTeX
  \usepackage{libertine}
}
\KOMAoption{captions}{tableheading}
\makeatletter
\@ifpackageloaded{caption}{}{\usepackage{caption}}
\AtBeginDocument{%
\ifdefined\contentsname
  \renewcommand*\contentsname{Содержание}
\else
  \newcommand\contentsname{Содержание}
\fi
\ifdefined\listfigurename
  \renewcommand*\listfigurename{Список иллюстраций}
\else
  \newcommand\listfigurename{Список иллюстраций}
\fi
\ifdefined\listtablename
  \renewcommand*\listtablename{Список таблиц}
\else
  \newcommand\listtablename{Список таблиц}
\fi
\ifdefined\figurename
  \renewcommand*\figurename{Рисунок}
\else
  \newcommand\figurename{Рисунок}
\fi
\ifdefined\tablename
  \renewcommand*\tablename{Таблица}
\else
  \newcommand\tablename{Таблица}
\fi
}
\@ifpackageloaded{float}{}{\usepackage{float}}
\floatstyle{ruled}
\@ifundefined{c@chapter}{\newfloat{codelisting}{h}{lop}}{\newfloat{codelisting}{h}{lop}[chapter]}
\floatname{codelisting}{Список}
\newcommand*\listoflistings{\listof{codelisting}{Листинги}}
\makeatother
\makeatletter
\makeatother
\makeatletter
\@ifpackageloaded{caption}{}{\usepackage{caption}}
\@ifpackageloaded{subcaption}{}{\usepackage{subcaption}}
\makeatother
\ifLuaTeX
\usepackage[bidi=basic]{babel}
\else
\usepackage[bidi=default]{babel}
\fi
\babelprovide[main,import]{russian}
\babelprovide[import]{english}
% get rid of language-specific shorthands (see #6817):
\let\LanguageShortHands\languageshorthands
\def\languageshorthands#1{}
\ifLuaTeX
  \usepackage{selnolig}  % disable illegal ligatures
\fi
\usepackage[style=gost-numeric,backend=biber,langhook=extras,autolang=other*]{biblatex}
\addbibresource{bib/cite.bib}
\usepackage{csquotes}
\usepackage{bookmark}

\IfFileExists{xurl.sty}{\usepackage{xurl}}{} % add URL line breaks if available
\urlstyle{same} % disable monospaced font for URLs
\hypersetup{
  pdftitle={Отчёта по лабораторной работе 5},
  pdfauthor={Нджову Нелиа},
  pdflang={ru-RU},
  hidelinks,
  pdfcreator={LaTeX via pandoc}}

\title{Отчёта по лабораторной работе 5}
\usepackage{etoolbox}
\makeatletter
\providecommand{\subtitle}[1]{% add subtitle to \maketitle
  \apptocmd{\@title}{\par {\large #1 \par}}{}{}
}
\makeatother
\subtitle{Математическое моделирование}
\author{Нджову Нелиа}
\date{}

\begin{document}
\maketitle

\renewcommand*\contentsname{Содержание}
{
\setcounter{tocdepth}{1}
\tableofcontents
}
\listoffigures
\listoftables
\setstretch{1.5}
\chapter{Цель
работы}\label{ux446ux435ux43bux44c-ux440ux430ux431ux43eux442ux44b}

Целью данной лабораторной работы является построение математической
модели хищник-жертва.

\chapter{Задание}\label{ux437ux430ux434ux430ux43dux438ux435}

\section{Вариант 64}\label{ux432ux430ux440ux438ux430ux43dux442-64}

Для модели «хищник-жертва»:

\[
\begin{cases}
\frac{dx}{dt} = -0.48x(t) + 0.031x(t)y(t)\\
\frac{dy}{dt} = 0.68y(t) - 0.031x(t)y(t)
\end{cases}
\](2.1)

Постройте график зависимости численности хищников от численности жертв,
а также графики изменения численности хищников и численности жертв при
следующих начальных условиях: \(x_0 = 3\), \(y_0 = 18\). Найдите
стационарное состояние системы.

\chapter{Теоретическое
введение}\label{ux442ux435ux43eux440ux435ux442ux438ux447ux435ux441ux43aux43eux435-ux432ux432ux435ux434ux435ux43dux438ux435}

Простейшая модель взаимодействия двух видов типа «хищник --- жертва»
модель Лотки-Вольтерры. Данная двувидовая модель основывается на
следующих предположениях:

\begin{enumerate}
\def\labelenumi{\arabic{enumi}.}
\item
  Численность популяции жертв x и хищников y зависят только от времени
  (модель не учитывает пространственное распределение популяции на
  занимаемой территории)
\item
  В отсутствии взаимодействия численность видов изменяется по модели
  Мальтуса, при этом число жертв увеличивается, а число хищников падает
\item
  Естественная смертность жертвы и естественная рождаемость хищника
  считаются несущественными
\item
  Эффект насыщения численности обеих популяций не учитывается
\item
  Скорость роста численности жертв уменьшается пропорционально
  численности хищников
\end{enumerate}

\[
\begin{cases}
\frac{dx}{dt} = ax(t) - bx(t)y(t)\\
\frac{dy}{dt} = -cy(t) + dx(t)y(t)
\end{cases}
\](3.1)

В этой модели x -- число жертв, y - число хищников. Коэффициент a
описывает скорость естественного прироста числа жертв в отсутствие
хищников, с естественное вымирание хищников, лишенных пищи в виде жертв.
Вероятность взаимодействия жертвы и хищника считается пропорциональной
как количеству жертв, так и числу самих хищников (xy). Каждый акт
взаимодействия уменьшает популяцию жертв, но способствует увеличению
популяции хищников (члены -bxy и dxy в правой части уравнения)

Найдём стационарное состояние системы (3.1). Для этого приравняем её
правые части к нулю

\[
\begin{cases}
ax(t)-bx(t)y(t) = 0\\
-cy(t)+dx(t)y(t) = 0
\end{cases}
\](3.2)

Из полученной системы (3.2) получаем, что стационарное состояние системы
будет в точке \(x_0 = c/d\), \(y_0 = a/b\). Если начальные значения
задать в стационарном состоянии \(x(0) = x_0\), \(y(0) = y_0\), то в
любой момент времени численность популяций изменяться не будет. При
малом отклонении от положения равновесия численности как хищника, так и
жертвы с течением времени не возвращаются к равновесным значениям, а
совершают периодические колебания вокруг стационарной точки.

\chapter{Выполнение лабораторной
работы}\label{ux432ux44bux43fux43eux43bux43dux435ux43dux438ux435-ux43bux430ux431ux43eux440ux430ux442ux43eux440ux43dux43eux439-ux440ux430ux431ux43eux442ux44b}

\section{Поиск стационарного состояния
системы}\label{ux43fux43eux438ux441ux43a-ux441ux442ux430ux446ux438ux43eux43dux430ux440ux43dux43eux433ux43e-ux441ux43eux441ux442ux43eux44fux43dux438ux44f-ux441ux438ux441ux442ux435ux43cux44b}

Найдём стационарное состояние системы (2.1). Для этого приравняем её
правые части к нулю

\[
\begin{cases}
  -0.48x(t)+0.031x(t)y(t) = 0\\
  0.68y(t)- 0.031x(t)y(t) = 0
\end{cases}
\](4.1)

\[
\begin{cases}
  -0.48 + 0.031y(t) = 0\\
  0.68 - 0.031x(t) = 0
\end{cases}
\](4.2)

Из полученной системы получаем, что стационарное состояние системы будет
в точке \(x_0 = 0.68/0.031 = 21.9355\), \(y_0 = 0.48/0.031 = 15.4839\).
Если начальные значения задать в стационарном состоянии \(x(0) = x_0\),
\(y(0) = y_0\), то в любой момент времени численность популяций
изменяться не будет. При малом отклонении от положения равновесия
численности как хищника, так и жертвы с течением времени не возвращаются
к равновесным значениям, а совершают периодические колебания вокруг
стационарной точки.

\section{Реализация в
Julia}\label{ux440ux435ux430ux43bux438ux437ux430ux446ux438ux44f-ux432-julia}

Зададим начальные значения функций и параметры, а также интервал
интегрирования(рис.1)

\begin{figure}

\centering{

\includegraphics[width=0.7\textwidth,height=\textheight]{image/pic10.png}

}

\caption{\label{fig-001}Реализация в Julia}

\end{figure}%

Зададим функцию для решения модели хищник-жертва. Для задания проблемы
используется функция \texttt{ODEProblem}, а для решения -- численный
метод \texttt{Tsit5()}(рис.2)

\begin{figure}

\centering{

\includegraphics[width=0.7\textwidth,height=\textheight]{image/pic11.png}

}

\caption{\label{fig-002}Реализация в Julia}

\end{figure}%

С помощью \texttt{plot} строим графики решения модели, а также фазовый
портрет (рис.3, 4 и 5)

\begin{figure}

\centering{

\includegraphics[width=0.7\textwidth,height=\textheight]{image/pic12.png}

}

\caption{\label{fig-003}Реализация в Julia}

\end{figure}%

\begin{figure}

{\centering \includegraphics[width=0.7\textwidth,height=\textheight]{image/pic2.png}

}

\caption{Решение модели при \(x_0 = 3, \, y_0=18\). Julia}

\end{figure}%

\includegraphics{image/pic3.png}\{\#fig:005width=70\%\}

Проверим найденное стационарное состояние с помощью вычислений в
программе, для этого используем деление соответствующих параметров и
получим \texttt{x\_c}, \texttt{y\_c}. Если это стационарное состояние
системы,то на графике решения мы не увидим изменения функций, а на
фазовом портрете получим одну точку. Поэтому зададим проблему со
стационарным состоянием в виде начального условия и решим ее(рис.6 и 7)

\begin{figure}

\centering{

\includegraphics[width=0.7\textwidth,height=\textheight]{image/pic13.png}

}

\caption{\label{fig-006}Реализация в Julia}

\end{figure}%

\begin{figure}

\centering{

\includegraphics[width=0.7\textwidth,height=\textheight]{image/pic1.png}

}

\caption{\label{fig-007}стационарной точки}

\end{figure}%

С помощью \texttt{plot} построим графики решения и фазовый портрет
(рис.8, 9 и 10)

\begin{figure}

\centering{

\includegraphics[width=0.7\textwidth,height=\textheight]{image/pic14.png}

}

\caption{\label{fig-008}Реализация в Julia}

\end{figure}%

\begin{figure}

{\centering \includegraphics[width=0.7\textwidth,height=\textheight]{image/pic4.png}

}

\caption{Решение модели при \(x_0 = x_c, \, y_0=y_c\). Julia}

\end{figure}%%
\begin{figure}

{\centering \includegraphics[width=0.7\textwidth,height=\textheight]{image/pic5.png}

}

\caption{Фазовый портрет модели при \(x_0 = x_c, \, y_0=y_c\). Julia}

\end{figure}%

Видим, что значения решений остаются неизменными, значит, мы верно
определили стационарное состояние системы.

\section{Реализация в
OpenModelica}\label{ux440ux435ux430ux43bux438ux437ux430ux446ux438ux44f-ux432-openmodelica}

Сначала зададим модель с данными в условиях варианта начальными
значениями.

\begin{verbatim}
// Параметры
a = 0.48;
b = 0.031;
c = 0.68;
d = 0.031;

// Система ДУ
function dx = syst2(t, x)
    dx(1) = -a*x(1) + b*x(1)*x(2);
    dx(2) = c*x(2) - d*x(1)*x(2);
endfunction

// Начальные условия
t0 = 0;
x0 = [3; 18];
t = [0:0.1:50];

// Решение
y = ode(x0, t0, t, syst2);

y1 = y(1, :);
y2 = y(2, :);

// ГРАФИК
scf(0);
clf();
plot(t, y1, 'b-', 'LineWidth', 2);
plot(t, y2, 'r--', 'LineWidth', 2);
xlabel('Время');
ylabel('Численность популяции');
title('Модель хищник-жертва (x_0=3, y_0=18)');
legend('Жертвы', 'Хищники');
xgrid();
savefig('lab5_2.png');

scf(1);
clf();
plot(y1, y2, 'k-', 'LineWidth', 2);
xlabel('x, жертвы');
ylabel('y, хищники');
title('Фазовый портрет');
xgrid();
savefig('lab5_2.1.png');
\end{verbatim}

После установки симуляции получим следующие графики решения модели и
фазового портрета (рис.11 и 12)

\begin{figure}

{\centering \includegraphics[width=0.7\textwidth,height=\textheight]{image/pic6.png}

}

\caption{Решение модели при \(x_0 = 3, \, y_0=18\). OpenModelica}

\end{figure}%%
\begin{figure}

{\centering \includegraphics[width=0.7\textwidth,height=\textheight]{image/pic7.png}

}

\caption{Фазовый портрет модели при \(x_0 = 3, \, y_0=18\).
OpenModelica}

\end{figure}%

Далее поменяем начальное значение на стационарное состояние системы:

\begin{verbatim}
// Стационарная точка
x_c = c / d;   // 0.68 / 0.031 ≈ 21.9355
y_c = a / b;   // 0.48 / 0.031 ≈ 15.4839

// Решение со стационарными начальными условиями
x0_stat = [x_c; y_c];
y_stat = ode(x0_stat, t0, t, syst2);

y1_stat = y_stat(1, :);
y2_stat = y_stat(2, :);

// ГРАФИК
scf(2);
clf();
plot(t, y1_stat, 'b-', 'LineWidth', 2);
plot(t, y2_stat, 'r--', 'LineWidth', 2);
xlabel('Время');
ylabel('Численность популяции');
title('Стационарное состояние: x0 = x_c, y0 = y_c');
legend('Жертвы', 'Хищники');
xgrid();
savefig('lab5_2.2.png');

scf(3);
clf();
plot(y1_stat, y2_stat, 'ko', 'MarkerSize', 10, 'MarkerFaceColor', 'k');
xlabel('x, жертвы');
ylabel('y, хищники');
title('Фазовый портрет (стационарная точка)');
xgrid();
savefig('lab5_2.3.png');
\end{verbatim}

После установки симуляции получим следующие графики решения модели и
фазового портрета (рис.13 и 14)

\begin{figure}

{\centering \includegraphics[width=0.7\textwidth,height=\textheight]{image/pic8.png}

}

\caption{Решение модели при \(x_0 = x_c, \, y_0=y_c\). OpenModelica}

\end{figure}%%
\begin{figure}

{\centering \includegraphics[width=0.7\textwidth,height=\textheight]{image/pic9.png}

}

\caption{Фазовый портрет модели при \(x_0 = x_c, \, y_0=y_c\).
OpenModelica}

\end{figure}%

Графики, полученнные в OpenModelica идентичны графикам в Julia.

\chapter{Выводы}\label{ux432ux44bux432ux43eux434ux44b}

Выполнив эту лабораторную работу, я построила математическую модель
хищник жертва и провели анализ.


\printbibliography


\end{document}
